%!TEX program = lualatex
\documentclass[10pt,journal,compsoc]{IEEEtran}

\usepackage{ifluatex}
\ifCLASSOPTIONcompsoc
  \usepackage[nocompress]{cite}
\else
  \usepackage{cite}
\fi

\ifCLASSINFOpdf
  \usepackage[pdftex]{graphicx}
\else
  \usepackage{graphicx}
\fi

\usepackage{array}

\usepackage{siunitx}
\usepackage{adjustbox}


\ifCLASSOPTIONcompsoc
  \usepackage[caption=false,font=footnotesize,labelfont=sf,textfont=sf]{subfig}
\else
  \usepackage[caption=false,font=footnotesize]{subfig}
\fi


\usepackage{stfloats}

\usepackage{url}

\usepackage{comment}
\usepackage{balance}

%\newcommand{\ada}[1]{\textcolor{red}{#1}}
\newcommand{\ada}[1]{}

\usepackage{tikz}
\usepackage[utf8]{inputenc}
\ifluatex
\usepackage{fontspec} % This line only for XeLaTeX and LuaLaTeX
\fi
\usepackage{pgfplots}
\pgfplotsset{compat=1.16}
\usepackage{color, colortbl}
\usetikzlibrary{svg.path}
\usepackage{scalerel}

\definecolor{orcidlogocol}{HTML}{A6CE39}
\tikzset{
  orcidlogo/.pic={
    \fill[orcidlogocol] svg{M256,128c0,70.7-57.3,128-128,128C57.3,256,0,198.7,0,128C0,57.3,57.3,0,128,0C198.7,0,256,57.3,256,128z};
    \fill[white] svg{M86.3,186.2H70.9V79.1h15.4v48.4V186.2z}
                 svg{M108.9,79.1h41.6c39.6,0,57,28.3,57,53.6c0,27.5-21.5,53.6-56.8,53.6h-41.8V79.1z M124.3,172.4h24.5c34.9,0,42.9-26.5,42.9-39.7c0-21.5-13.7-39.7-43.7-39.7h-23.7V172.4z}
                 svg{M88.7,56.8c0,5.5-4.5,10.1-10.1,10.1c-5.6,0-10.1-4.6-10.1-10.1c0-5.6,4.5-10.1,10.1-10.1C84.2,46.7,88.7,51.3,88.7,56.8z};
  }
}

\newcommand\orcidicon[1]{\href{https://orcid.org/#1}{\mbox{\scalerel*{
\begin{tikzpicture}[yscale=-1,transform shape]
\pic{orcidlogo};
\end{tikzpicture}
}{|}}}}

\hyphenation{op-tical net-works semi-conduc-tor}

\usepackage[bookmarks=false]{hyperref} %<-- load after all other packages

\begin{document}

\title{Unexpected Performance of Intel\textsuperscript{\textregistered} Optane\textsuperscript{\small TM}
DC Persistent Memory}

\author{Tony~Mason~\orcidicon{0000-0002-0651-5019},~\IEEEmembership{Student~Member,~IEEE,}
Thaleia~Dimitra~Doudali~\orcidicon{0000-0002-3197-839X},~\IEEEmembership{Student~Member,~IEEE}, \\
Margo~Seltzer~\orcidicon{0000-0002-2165-4658},
        and~Ada~Gavrilovska~\orcidicon{0000-0003-4199-2512},~\IEEEmembership{Member,~IEEE}% <-this % stops a space
\thanks{T. Mason and M. Seltzer are with the Department of Computer Science, University of British Columbia, Vancouver, BC Canada e-mail: fsgeek@cs.ubc.ca, mseltzer@cs.ubc.ca.}% <-this % stops a space
\thanks{T. Mason, T. Doudali, and A. Gavrilovska are with the College of Computing, Georgia Institute of Technology, Atlanta,
GA, 30332 USA e-mail: fsgeek@gatech.edu, thdoudali@gatech.edu, ada@cc.gatech.edu.}% <-this % stops a space
\thanks{Manuscript received August 13, 2019; revised \today.}
}

\markboth{IEEE COMPUTER ARCHITECTURE LETTERS}%
{Mason \MakeLowercase{\textit{et al.} Unexpected Performance of Intel\textsuperscript{\textregistered} Optane\textsuperscript{\tiny TM}
DC Persistent Memory}}

\IEEEtitleabstractindextext{%
\begin{abstract}
  %%
  We evaluated Intel\textsuperscript{\textregistered} Optane\textsuperscript{\tiny TM} DC Persistent Memory and found that Intel's persistent memory is highly sensitive to data locality, size, and access patterns, which becomes clearer by optimizing both virtual memory page size and data layout for locality. Using the \textit{Polybench} high-performance computing benchmark suite and controlling for mapped page size, we evaluate PMEM performance relative to DRAM. In particular, the Linux persistent memory (PMEM) support automatically maps persistent memory in large pages.  The effect of large pages for PMEM and small pages for DRAM is enormous, dwarfing other effects discussed in the literature.  We found PMEM performance comparable to DRAM performance for a number of tests with significant performance improvements when optimized for data locality.
  %%  
\end{abstract}

% Note that keywords are not normally used for peerreview papers.
\begin{IEEEkeywords}
Persistent Memory
\end{IEEEkeywords}
}%


\maketitle


\IEEEdisplaynontitleabstractindextext
\IEEEpeerreviewmaketitle

\IEEEraisesectionheading{\section{Introduction}\label{sec:introduction}}

\balance

%%
\IEEEPARstart{N}{on-volatile} memory research began decades ago~\cite{wu1994envy}.
The April 2019 release of Intel Optane\textsuperscript{\tiny TM} DC Persistent Memory (PMEM) has accelerated
that research. Prior to availability, most work was done via emulation~\cite{wu1994envy,Maciejewski2017persistent,dulloor2014system,volos2015quartz,doudali2017comerge}.  We expected PMEM to exhibit lower bandwidth and higher latency compared to DRAM; work with PMEM is consistent with these expectations~\cite{gill2019single,izraelevitz2019basic}.  Yet we also note the actual performance behavior is more complex~\cite{peng2019system}. 
%%

%%
We evaluated PMEM using \textit{Polybench}, a suite of well-known microbenchmarks designed to measure the efficiency of memory locality optimization techniques.  Our initial results, shown in Figure \ref{fig:polybench}, defied the findings of our prior emulation study~\cite{doudali2017comerge}.  We investigated this behavior to better understand the performance profile of PMEM. We confirmed that locality is critical and identified multiple ways in which locality manifests itself in current PMEM systems.  We verified our observations on multiple different systems.  We report our results on the last system that we evaluated.
%%

%%
We observe that the most significant effect of locality is related to default memory management policy, with secondary effects from the PMEM memory controller's use of striping, read-ahead, and write-behind caching.  The Linux default memory management policy is to use the largest possible pages for PMEM, while it defaults to 4KB pages for DRAM.  DAX-aware file systems (for application sharing, dynamic allocation, and security) interact with this policy as well.  We observed that the impact of this is up to 5x and we conclude that those publishing  PMEM evaluations should control for page size. PMEM striping primarily benefits workloads with good data locality (within the stripe size) and high bandwidth requirements.  PMEM without striping often performs similarly when high bandwidth is less critical.  PMEM caching in the memory controller and memory modules benefits CPUs with smaller caches.  Once we controlled for page size and used a memory locality optimizing tool, we confirmed performance for PMEM was as good or better than expectations.  We achieved this performance from careful optimization, \textit{not} by using default configurations.
%%

\begin{figure}[!htb]
    \captionsetup{justification=centering}
    \centering
    \caption{Polybench: Performance of Striped PMEM relative to DRAM, Linux 5.0 kernel.  Surprising results showing better performance on PMEM are highlighted in \textbf{bold}. Sorted by PMEM performance.}
    \vspace{0.1cm}
    \label{fig:polybench}
    \begin{tabular}{c}
        \includegraphics[width=0.495\textwidth]{pbfig-dec2019.eps}
    \end{tabular}
\end{figure}



\section{Background}\label{sec:methodology}

\begin{figure}[!th]
  \captionsetup{justification=centering}
  \centering
  \caption{Polybench: DRAM vs Interleaved PMEM (both 4K pages).  Results closer to expectations but still better than expected.  Caching masks lower bandwidth of PMEM; memory usage optimization improves cache benefits which improves performance and decreases benefit of PMEM interleaving.}
  \vspace{2mm}
  \label{fig:polybench:dram-vs-ipmem-4k}
  \begin{tabular}{c}
    \includegraphics[width=0.45\textwidth]{pb-4k-DRAM-4K-iPMEM.eps}
  \end{tabular}
\end{figure}


%%
PMEM's AppDirect mode permits two types of direct-access (DAX) usage: \textit{devdax} access, which provides raw PMEM that is memory mapped by an application for exclusive use. Devdax mode requires pre-selecting memory allocation units (4KB, 2MB, and 1GB), static partitioning of the overall memory, using the ``ndctl'' utility, and is restricted to privileged applications; and \textit{fsdax} access, which uses a DAX-aware file system to provide support for sharing between multiple applications, dynamic memory allocation --- including page alignment and allocation unit size, and multi-user security.  When an application memory maps a file on a DAX-aware file system the operating system configures the process page tables to directly reference the physical addresses of the underlying PMEM.  Preferred practice is to use a DAX-aware file system~\cite{rudoff2019NVMProgrammingModel}.
%%


%%
The Persistent Memory Development Kit (PMDK) includes libraries for transparently converting standard memory allocation calls (e.g., \textit{malloc}) to a memory allocator using arenas backed directly by files on a DAX filesystem device.  This is similar to the approach used by \textit{hugetlbfs}, a Linux file sytem that exposes DRAM to applications for allocation against large pages --- the standard malloc calls in that case are implemented using arenas backed directly by files on hugetlbfs.  We use both of these mechanisms in our evaluation of the Polybench tests.
%%

%%
Our PMEM hardware uses two dedicated memory controllers per CPU; each memory controller has three channels; each channel can manage two PMEM modules, which are equivalent to DRAM DIMM packages.   The memory controllers include a small cache memory and provide the ability to transparently stripe across the PMEM modules although non-interleaved access is also supported.  Each PMEM module also contains memory for converting 64 byte cache line sized load/store operations into 256 byte PMEM block size load/store operations.
%%

%%
The Linux operating system includes native support for PMEM.  However, the policies for PMEM differ from the ones for DRAM. The Linux kernel transparently uses large memory page mappings for PMEM whenever possible, as dictated by the alignment and length of the page mappings.  Linux version 5.3 uses 1GB, 2MB, or 4KB page sizes for PMEM when setting up page tables for running processes.  Earlier Linux versions only used 2MB or 4KB pages for PMEM. The default for DRAM remains 4KB.  We note that large page performance impact is well-described in the literature~\cite{agarwal2017thermostat,panwar2018making}.  While Linux behavior can be tuned, because default page sizes differ for PMEM and DRAM, it can create significant performance differences --- we observed up to 5x difference.
%%


%%
We evaluated three different DAX-aware file systems: ext and xfs, which are standard Linux file systems incorporating experimental DAX support, and NOVA. The NOVA file system was the first purpose-built file system for PMEM~\cite{xu2016nova}.  However, we found NOVA unsuitable for use as a DAX file system, because its storage allocation policy does not preserve the 2MB page alignment the OS requires to provide large page support; this is not reported in even recent literature~\cite{izraelevitz2019basic,yang2019empirical}. Initially, we had not expected our choice of DAX-aware file systems to be relevant, as once memory mapped, the PMEM is directly accessed by the application.  As a result of our discovery, our evaluation in this paper uses ext4, which did not exhibit the behavior observed with NOVA.
%%

\begin{figure}[!b]
  \captionsetup{justification=centering}
  \centering
  \caption{Polybench: Cases where PMEM (non-interleaved PMEM) faster than iPMEM (interleaved PMEM) --- interleaving does \textit{not} provide substantial benefit in these cases; surprising because we expect the 2MB page size to dominate.}
  \vspace{2mm}
  \label{fig:polybench:PMEM-vs-ipmem}
  \begin{tabular}{c}
    \includegraphics[width=0.45\textwidth]{pb-2M-PMEM-2MB-iPMEM-non-interleaved-faster.eps}
  \end{tabular}
\end{figure}



%%
The Polybench test suite is a set of 30 computational kernels drawn from several different domains and is commonly used for evaluating the impact of memory performance with respect to those tasks. We previously used Polybench~\cite{doudali2017comerge}, and Polybench continues to be actively used for evaluating memory optimizer performance.
%%



% Need to define this somewhere
\definecolor{Gray}{gray}{0.9}
\begin{table*}[!ht]
	\caption{Polybench Test Results: data size, by memory type (DRAM, PMEM, iPMEM), page size (4K, 2M), and optimization (\texttt{-O3} versus \texttt{-poly}).  Execution time in seconds; sorted by 4K DRAM time improvement (optimized versus poly). Only shows test where polyhedral optimization benefitted at least one memory configuration (before rounding); other results omitted. Memory locality impact differs across memory types and configurations. CPU and memory allocation bound to one NUMA node. Ratio > 1 (pre-rounding) shows polyhedral optimization benefits.}
	\label{table:poly2}
  \begin{adjustbox}{max width=\textwidth}
    \begin{tabular}{r|S|SSS|SSS|SSS|SSS|SSS|SSS}
      \multicolumn{1}{c|}{Test} &
      \multicolumn{1}{c|}{Data} &
      \multicolumn{6}{c|}{DRAM} &
      \multicolumn{6}{c|}{PMEM} &
      \multicolumn{6}{c}{iPMEM}
      \tabularnewline
      & 
      \multicolumn{1}{c|}{Size} & 
      \multicolumn{3}{c}{4K} &
      \multicolumn{3}{c|}{2M} &
      \multicolumn{3}{c}{4K} &
      \multicolumn{3}{c|}{2M} &
      \multicolumn{3}{c}{4K} &
      \multicolumn{3}{c}{2M}
      \tabularnewline
        & {MB} 
        & {Opt}  & {Poly} & {Ratio} & {Opt} & {Poly} & {Ratio} 
        & {Opt}  & {Poly} & {Ratio} & {Opt} & {Poly} & {Ratio} 
        & {Opt}  & {Poly} & {Ratio} & {Opt} & {Poly} & {Ratio}
      \tabularnewline
      \hline
      gemm & 91.6 
      & 29.3 & \textbf{1.5} & 19.1 & 14.1 & \textbf{1.5} &  9.3 
      & 39.0 & \textbf{1.7} & 22.6 & 32.0 & \textbf{1.7} & 18.7 
      & 37.8 & \textbf{1.7} & 22.9 & 27.5 & \textbf{1.7} & 16.6
      \tabularnewline
      \rowcolor{Gray}
      3mm & 120.2 
      & 17.1 & \textbf{1.5} & 11.8 & 12.1 & \textbf{1.4} & 8.4
      & 16.6 & \textbf{1.8} &  9.5 & 11.9 & \textbf{1.7} & 6.8
      & 17.1 & \textbf{1.7} & 10.3 & 11.7 & \textbf{1.6} & 7.2
      \tabularnewline
      2mm & 85.8 
      & 5.4 & \textbf{0.5} & 11.3 &  3.8 & \textbf{0.5} & 7.7
      & 5.2 & \textbf{0.6} &  8.6 &  4.1 & \textbf{0.6} & 6.9
      & 5.4 & \textbf{0.6} &  9.7 &  3.9 & \textbf{0.6} & 6.9
      \tabularnewline
      \rowcolor{Gray}
      correlation & 137.4 
      & 105.1 & \textbf{22.4} & 4.7 & 23.2 & \textbf{20.5} & 1.1 
      & 123.8 & \textbf{23.9} & 5.2 & 45.5 & \textbf{23.3} & 2.0 
      & 118.8 & \textbf{23.6} & 5.0 & 39.9 & \textbf{22.9} & 1.7
      \tabularnewline
      covariance & 49.5 
      & 19.7 & \textbf{4.7} & 4.2 & 3.3 & 4.4 & 0.8
      & 20.3 & \textbf{5.0} & 4.1 & 4.4 & 4.9 & 0.9
      & 20.0 & \textbf{4.9} & 4.1 & 4.1 & 4.8 & 0.9
      \tabularnewline
      \rowcolor{Gray}
      mvt & 1717.1 
      & 3.4 & \textbf{1.0} & 3.5 & 1.1 & \textbf{1.0} & 1.0
      & 4.7 & \textbf{2.0} & 2.3 & 2.3 & \textbf{2.0} & 1.1
      & 4.5 & \textbf{1.9} & 2.4 & 2.1 & \textbf{1.9} & 1.1
      \tabularnewline
      gemver & 1717.4 
      & 3.7 & \textbf{1.5} & 2.4 & 1.2 & 1.7 & 0.8
      & 7.9 & \textbf{4.9} & 1.6 & 5.3 & \textbf{4.9} & 1.1
      & 5.2 & \textbf{3.2} & 1.6 & 2.7 & 3.2 & 0.9
      \tabularnewline
      \rowcolor{Gray}
      jacobi-1d & 152.6 
      & 27.6 & \textbf{12.0} & 2.3 & 26.0 & \textbf{11.0} & 2.4
      & 261.8 & \textbf{70.0} & 3.7 & 260.0 & \textbf{70.1} & 3.7
      & 89.0 & \textbf{70.0} & 1.3 & 88.6 & \textbf{69.9} & 1.3
      \tabularnewline
      doitgen & 512.5 
      & 13.0 & \textbf{6.2} & 2.1 & 8.6 & \textbf{6.1} & 1.4
      & 13.3 & \textbf{10.2} & 1.3 & 12.4 & \textbf{9.8} & 1.3
      & 13.5 & \textbf{7.2} & 1.4 & 12.5 & \textbf{7.1} & 1.8
      \tabularnewline
      \rowcolor{Gray}
      gramschmidt & 91.6 
      & 67.0 & \textbf{38.8} & 1.7 & 11.3 & \textbf{9.75} & 1.2
      & 73.1 & \textbf{50.6} & 1.4 & 24.5 & 28.7 & 0.9
      & 70.5 & \textbf{44.8} & 1.6 & 16.6 & 19.1 & 0.9
      \tabularnewline
      fdtd-2d & 1464.8 
      & 33.0 & \textbf{19.3}  & 1.7 & 34.8  & \textbf{17.0} & 2.0
      & 275.5 & \textbf{57.4} & 4.8 & 273.6 & \textbf{56.9} & 4.8 
      & 104.29 & \textbf{52.2} & 2.0 & 103.8 & \textbf{52.5} & 2.0
      \tabularnewline
      \rowcolor{Gray}
      syrk & 61.0 
      & 9.1 & \textbf{5.4}  & 1.6 & 9.0 & \textbf{5.4}  & 2.0
      & 15.1 & \textbf{6.0} & 2.5 & 15.0 & \textbf{5.9} & 2.5
      & 14.6 & \textbf{6.1} & 2.4 & 15.1 & \textbf{6.1} & 2.5
      \tabularnewline
      jacobi-2d & 976.6 
      & 18.7  & \textbf{12.1} & 1.5 & 18.1  & \textbf{9.0}  & 2.0 
      & 179.6 & \textbf{62.3} & 2.8 & 179.6 & \textbf{62.1} & 2.9
      & 59.6  & 60.1          & 1.0 & 61.0  & \textbf{60.0} & 1.0
      \tabularnewline
      \rowcolor{Gray}
      syr2k & 91.6 
      & 17.9 & \textbf{12.0} & 1.5 & 18.0 & \textbf{11.8} & 1.5
      & 31.0 & \textbf{14.6} & 2.1 & 30.0 & \textbf{14.6} & 2.1
      & 26.0 & \textbf{14.3} & 2.8 & 25.9 & \textbf{14.2} & 2.8
      \tabularnewline
      symm & 91.6 
      & 49.6 & \textbf{34.2} & 1.4 & 9.8 & 11.5 & 0.9
      & 65.1 & \textbf{44.5} & 1.5 & 41.4 & \textbf{38.9} & 1.1
      & 64.4 & \textbf{44.0} & 1.5 & 37.8 & \textbf{23.5} & 1.6
      \tabularnewline
      \rowcolor{Gray}
      dynprog & 245.4 
      & 11.3 & \textbf{11.3} & 1.0 & 11.5 & 11.5          & 1.0
      & 71.4 & 73.0          & 1.0 & 74.7 & 75.3          & 1.0
      & 31.4 & 31.5          & 1.0 & 32.0 & \textbf{31.8} & 1.0
      \tabularnewline
      durbin & 1526.2 
      & 1.7 & \textbf{1.7}  & 1.0 & 0.5 & \textbf{0.5} & 1.0
      & 4.0 & 4.0           & 1.0 & 2.8 & \textbf{2.8} & 1.0
      & 2.6 & \textbf{2.6}  & 1.0 & 1.3 & 1.3          & 1.0
      \tabularnewline
      \rowcolor{Gray}
      seidel-2d & 190.7 
      & 25.1 & \textbf{25.1} & 1.0 & 24.5 & \textbf{24.5} & 1.0
      & 31.3 & \textbf{31.2} & 1.0 & 31.0 & 31.1          & 1.0
      & 25.1 & \textbf{25.1} & 1.0 & 25.1 & \textbf{25.1} & 1.0
      \tabularnewline
      trisolv & 1716.8 
      & 0.2 & 0.2          & 1.0 & 0.2 & \textbf{0.2} & 1.0
      & 0.3 & \textbf{0.3} & 1.0 & 0.3 & \textbf{0.3} & 1.0 
      & 0.3 & 0.3          & 1.0 & 0.3 & \textbf{0.3} & 1.0
      \tabularnewline
      \rowcolor{Gray}
      cholesky & 122.1 
      & 12.0 & 13.0 & 0.9 & 11.9 & \textbf{10.8} & 1.1
      & 18.8 & 20.5 & 0.9 & 18.5 & 19.4          & 1.0
      & 17.7 & 20.4 & 0.9 & 17.2 & 18.5          & 0.9
      \tabularnewline
      lu & 122.1 
      & 17.0 & 25.2            & 0.7 & 16.0 & 28.4            & 0.6
      & 209.8 & \textbf{192.2} & 1.1 & 219.2 & \textbf{192.8} & 1.1
      & 55.1 & 78.2            & 0.7 & 54.4 & 80.5            & 0.7
      \tabularnewline
      \rowcolor{Gray}
      reg\_detect & 381.6 
      & 30.1 & 52.2 & 0.6 & 28.7 & 51.0 & 0.6
      & 331.2 & \textbf{119.0} & 2.8 & 327.7 & \textbf{118.9} & 2.8
      & 93.3 & 118.6 & 0.8 & 92.6 & 118.4 & 0.8
      \tabularnewline
    \end{tabular}
  \end{adjustbox}
\end{table*}


\section{Evaluation}\label{sec:evaluation}
\ada{I prefer it when the evaluation is presented in present tense
  (except for when mentioning prior results). I do like the past tense
in the abstract of this paper, so maybe it will be weird to switch to
present tense. }

\begin{figure*}[!th]
  \captionsetup{justification=centering}
  \centering
  \caption{Polyhedral Memory optimizations for interleaved PMEM, 4K and 2M pages.  Varying cache layers impact performance differently. Results sorted by execution time without polyhedral optimization (\texttt{-polly} option to clang); boldface for improved poly performance.}
  \vspace{0.1cm}
  \label{fig:polybench:all:time}
  \begin{tabular}{cc}
    \includegraphics[width=0.45\textwidth]{pb-2M-iPMEM-poly-2M-iPMEM.eps} &
    \includegraphics[width=0.45\textwidth]{pb-4K-iPMEM-poly-4K-iPMEM.eps}
  \end{tabular}
\end{figure*}


%\input{appendix/polybench/aep/figures/tex/covariance-figure.tex}
%\input{appendix/polybench/aep/figures/tex/symm-figure.tex}
%\input{appendix/polybench/aep/figures/tex/lu-figure.tex}
%\input{appendix/polybench/aep/figures/tex/gemm-figure.tex}

%%
% We used the following tools for our evaluation: \textit{Polybench} 3.2~\cite{polybench,yuki2014understanding,yuki2015polybench}, a suite of benchmarks commonly used to evaluate memory and processor intensive workloads across a range of domains, and LLVM 10.0 --- specifically the clang compiler and polly polyhedral optimizations.

%%
Our evaluation system is a dual socket NUMA architecture system, using two Intel\textsuperscript{\textregistered} second generation Xeon\textsuperscript{\textregistered} Scalable processors (codename \textit{Cascade Lake}) running between 1.0GHz and 3.9GHz (variable) with 1.3MB L1 caches, 40MB L2 caches, and a 55MB L3 cache, 20 cores per processor, and 12 memory slots per processor, with six 32GB DRAM modules and six 256GB PMEM modules, all running at 2666MT/s. Each processor has two persistent memory controllers, with three channels per controller.  Each channel managed two PMEM modules; Intel refers to this as the 2-2-2 configuration~\cite{intel2019quickstart}. We found similar results on the same base system with different model Intel CPUs: 1-3.7GHz (variable), 32KB L1, 1MB L2, 32MB L3 caches.
%%

%%
We used Fedora 31 with Linux kernel 5.0.0, which includes native PMEM support. We configured the PMEM as six-way interleaved memory for one processor (iPMEM), and six single non-interleaved memories for the other processor (PMEM).  We used the SNIA standard programming model, which is a DAX-aware file system providing dynamic PMEM management and security; while it is possible to use devdax mode, which provides raw PMEM access, it requires static allocation and privileged (``root'') access.  We also used the PMDK libraries for transparently converting standard malloc calls into corresponding mmap calls for direct access PMEM, which only works with a DAX-aware file system~\cite{PMDK}.
%%


%%
We used Polybench/C Version 3.2~\cite{polybench,yuki2014understanding,yuki2015polybench} compiled with clang version 10.0.0, choosing compile time constants to report execution time and the Linux real-time scheduler, and using maximum (\texttt{-O3}) optimization. We tested both with (\texttt{--llvm -polly}) and without polyhedral optimization. We choose our dataset sizes to match our prior work~\cite{doudali2017comerge}.  We ran the single-threaded Polybench 3.2 tests serially, bound to a single core, and all memory, including PMEM, was allocated from memory local to the NUMA node of that core.  We used the PMDK allocator to redirect standard memory allocation calls to use memory mapped PMEM, via the DAX aware file system.  The PMDK uses a jemalloc-based memory allocator.
%%

%%
Figure \ref{fig:polybench} reproduces our initial results, using the ext4 file system properly configured to use aligned 2MB pages.  These initial results surprised us because they showed better performance than we expected.  We determined that while the default behavior for Linux is to use 4KB TLB mappings for DRAM, it used 2MB TLB mappings for PMEM when possible. Indeed, this behavior of Linux has changed; as of Linux 5.3 it now uses the largest possible TLB mapping for PMEM between 1GB, 2MB, and 4KB, based upon the alignment and length of the memory region being accessed.  This behavior has not previously been reported in the literature, despite its substantial impact on performance.  We counted the number of TLB and last level cache misses when forcing different page sizes; it accounts for all of the performance difference. Recent work has alluded to the 2MB page impact, but fails to explain why this occurs~\cite{kadekodi2019splitfs}, while other recent work fails to address it, despite the profound impact it has on performance for larger workloads~\cite{dong2019performance,peng2019system}. We note that NOVA, which was designed for PMEM and is the \textit{de facto} standard for PMEM-optimized file systems interacts poorly with the page-size needs of applications converted to use PMEM, because it mixes 4KB and 2MB page allocations internally.  Over time this leads to fragmentation, which causes the performance degradation we observed.  Thus, we subsequently switched to using ext4, which did not exhibit the performance degradation.
%%

%%
We used \texttt{perf}, a standard Linux performance utility, for measuring processor performance.  From perf, we identified dramatically different data TLB miss rates, which in turn led us to finding both the default behavior of Linux and the sensitivity of page locality to DAX file system allocation policy.  \textbf{OS implementers make fundamental decisions we must understand to achieve good performance.} Further, these choices do change over time.
%%

%%
We used \texttt{AEPWatch}, an open source Intel utility~\cite{IntelAEPWatch}, for measuring the performance behavior of the persistent memory controllers. Once we controlled for page size, this provided us with greater insight into the impact of read caching, write combining, and striping in the persistent memory controllers.  We found that \textbf{no single configuration produced consistently best performance.}  We will explore this further in future work.
%%

%%
We hypothesized that software memory locality techniques would also yield improved performance.  To test this, we used the LLVM compiler tools, which include a polyhedral memory optimizer. The memory optimizer optimizes memory layout and code generation to improve data locality; Polybench was constructed to evaluate the effectiveness of polyhedral optimization.  In Table \ref{table:poly2} we show our results across DRAM, (non-interleaved) PMEM, and (interleaved) iPMEM, for both 4KB and 2MB page sizes.  Most PMEM configurations benefitted from polyhedral optimization, which reflects the sensitivity of the system to locality. We found that non-interleaved PMEM benefitted most from polyhedral optimization, often exhibiting comparable performance to interleaved PMEM.  In only two cases did PMEM exhibit better performance than DRAM, likely due to the additional caching in the PMEM system itself. We omit results for eight tests where no memory configuration benefitted from polyhedral optimization. We will explore how to improve the polyhedral optimizer for use with persistent memory in future work.
%%

%%
Figure \ref{fig:polybench:dram-vs-ipmem-4k} was a fairer comparison of memory performance and a clear indication on the importance of locality.  We evaluated the impact of polyhedral memory optimization to corroborate our theory on the importance of data locality.  Figure \ref{fig:polybench:PMEM-vs-ipmem} was surprising to us because it suggested that locality was a more important factor than memory striping. Figure \ref{fig:polybench:all:time} demonstrated that in many cases polyhedral memory optimization improved data locality by as much as 80\%.  Table \ref{table:poly2} provides the specific timings for the 22 tests where polyhedral optimization improved performance for at least one memory type.  However, we were unable to identify any single factor that explained these exact results: we considered both data and instruction TLB miss rates, last level cache misses, data set sizes, and instruction counts. We suspect this is due to the generic nature of the polyhedral optimizer, which has no specific knowledge of the processor or memory characteristics on our test system; it is also possible there are secondary features within the PMEM itself that, while not exposed, impact this.  We will explore providing additional guidance for optimization in future work.

\begin{comment}
%%
While prior work~\cite{yang2019empirical} combined these behaviors into a single measure (Effective Write Ration -- ``EWR'') our evaluation suggests that the performance profile is more nuanced than captured by this single ratio. Thus, we expect systems using PMEM will require workload specific evaluation across the range of configuration options to optimize performance.
%%
\end{comment}





\begin{comment}
\section{Future Work}

Our results suggest a number of areas for future work: (1) we need to evaluate the existing polyhedral optimizations to determine why they do not benefit $\approx \frac{1}{3}$ of the Polybench tests.  We expect we will use detailed study of PMEM behavior to better understand the impact on data locality.  (2) We anticipate better understanding of dataset partitioning across multiple caching layers with different optimization functions could yield better insights into developing memory optimization tools. (3) We theorize that memory optimizaiton tools should provide multiple potential implementations optimized to explit specific memory topologies.  (4) We propose developing a DAX-optimized file system that focuses on providing optimal performance of PMEM as well as application sharing, dynamic allocation, and security focused on maximizing performance of applications using PMEM via the DAX model.
\end{comment}

\section{Conclusion}

%%
We found PMEM is more sensitive to memory locality than DRAM; this is due to differences in the hardware implementing these memory technologies.  The decision to default PMEM to large pages provides a substantial boost in page locality.  As datasets grow in size, using large pages yields better TLB locality because more actual PMEM can be described with a fixed number of TLB entries.  PMEM's support for interleaving is most beneficial for data with both good locality and high bandwidth requirements; non-interleaved PMEM works comparably well to interleaved PMEM with good data locality and moderate bandwidth requirements, such as with CPUs employing larger caches.
%%

%%
We caution those evaluating PMEM to ensure the evaluation properly controls for page size because it's impact may lead to incorrect results.  Those using PMEM should consider optimizing data locality as  it can lead to substantial performance gains.  However, those gains are clearly sensitive to the behavior of the underlying memory and thus are dependent upon the workload. Not optimizing risks sacrificing performance when using PMEM.
%%

%%
While our evaluation was done with the first generation of commercially available PMEM, we expect these insights will generalize to future implementations of PMEM, particularly those related to the impact of data locality, OS behavior, and caching/striping (interleaving) of PMEM.
%%


\ifCLASSOPTIONcompsoc
  \section*{Acknowledgments}
\else
  \section*{Acknowledgment}
\fi

%%
We thank the anonymous reviewers for their valuable feedback, as well as Vaastav Anand, Swati Goswami, Nodir Kodirov, Joel Nider, Ranjan Sarpangala Venkatesh, and Brian Veraa for their assistance in reviewing drafts. This research was partially funded under NSF award 1822972, US Dept. of Energy UNITY, Intel Corporation, and Exascale Computing Project SICM (Simplified Interface to Complex Memory).
%%

\nocite{Hunter:2007} % this is the matplotlib reference
\nocite{yang2019empirical} % latest UCSD paper that includes EWR
\nocite{peng2019system} % LLNL MEMSYS paper
\nocite{yuki2014understanding} % Understanding Polybench
\nocite{yuki2015polybench} % Polybench 4.0
\nocite{xu2016nova} % NOVA File System, damn it all to hell

\bibliographystyle{IEEEtranS}
\bibliography{calpub}
\end{document}
